\section{Bases de Dise\~no}

El presente estudio se fundamenta en los datos de geometr\'ia del tanque, las propiedades termof\'isicas del fluido almacenado y las condiciones operacionales reportadas por Ingredion~S.A. A continuaci\'on se detallan las bases de dise\~no que sustentan los c\'alculos de transferencia de calor.

\subsection{Informaci\'on Suministrada por Ingredion~S.A.}

La informaci\'on de proceso y las especificaciones del fluido utilizadas en este estudio fueron entregadas por Ingredion~S.A. como parte del alcance del proyecto. La Figura~\ref{fig:datos_ingredion} reproduce el documento de registro suministrado por el cliente para el tanque de almacenamiento de glucosa Tag~53A-90A-0056.

\begin{figure}[H]
\centering
\includegraphics[width=0.75\textwidth]{../../datosingredion.png}
\caption{Informaci\'on suministrada por Ingredion~S.A. para el tanque de almacenamiento de glucosa Tag~53A-90A-0056.}
\label{fig:datos_ingredion}
\end{figure}

A partir de la informaci\'on anterior, las condiciones de dise\~no del equipo y del servicio de calentamiento se resumen en la Tabla~\ref{tab:datos_ingredion_equipo}. Estos datos constituyen el punto de partida para la definici\'on de los escenarios de an\'alisis t\'ermico.

\begin{table}[H]
\centering
\caption{Condiciones del tanque y del servicio de calentamiento seg\'un informaci\'on entregada por Ingredion~S.A.}
\label{tab:datos_ingredion_equipo}
\setcounter{itemcount}{0}
\begin{tabularx}{\textwidth}{@{}NXclX@{}}
\toprule
\multicolumn{1}{c}{\textbf{\'Item}} & \textbf{Par\'ametro} & \textbf{Valor} & \textbf{Unidad} & \textbf{Observaci\'on} \\
\midrule
& Tag del tanque & 53A-90A-0056 & --- & Tanque de almacenamiento de glucosa \\
& Fluido almacenado & Glucosa Globe 1130 & --- & Jarabe de ma\'iz grado alimenticio \\
& Temperatura objetivo de almacenamiento en el fondo & 60 & \textdegree C & Condici\'on de operaci\'on \\
& Material del fondo & SS316L & --- & Acero inoxidable \\
& Espesor del fondo & 9,0 & mm & Seg\'un planos del fabricante \\
& Frecuencia de descarga & 5 & cargas/d\'ia & Requerimiento operativo \\
& Volumen por descarga & 23 -- 24 & m$^3$ & Cada evento de descarga \\
& Temperatura del agua de calentamiento & 65 & \textdegree C & Servicio industrial disponible \\
& Caudal de agua de calentamiento & 20,3 & m$^3$/h & Informaci\'on de proceso \\
& Presi\'on del agua de calentamiento & 30 & PSI & Informaci\'on de proceso \\
& Di\'ametro de tuber\'ia de alimentaci\'on & 2 & pulgadas & SS~304, Sch~40 \\
\bottomrule
\end{tabularx}
\end{table}

La Tabla~\ref{tab:datos_ingredion_especificaciones} presenta las especificaciones f\'isico-qu\'imicas del jarabe Globe 1130 reportadas por Ingredion~S.A., las cuales se emplearon para la caracterizaci\'on termof\'isica del fluido almacenado.

\begin{table}[H]
\centering
\caption{Especificaciones f\'isico-qu\'imicas de la glucosa Globe 1130 seg\'un informaci\'on entregada por Ingredion~S.A.}
\label{tab:datos_ingredion_especificaciones}
\setcounter{itemcount}{0}
\begin{tabularx}{\textwidth}{@{}NlccX@{}}
\toprule
\multicolumn{1}{c}{\textbf{\'Item}} & \textbf{Propiedad} & \textbf{M\'in} & \textbf{M\'ax} & \textbf{Observaci\'on} \\
\midrule
& Densidad aparente & 1,43 & --- & kg/L a condiciones de referencia \\
& Sustancia seca (\textdegree Brix) & 80,5 & 81,6 & Base operativa del estudio: 80,6~\textdegree Brix \\
& Equivalente de dextrosa (DE) & 38 & 42 & Conversi\'on regular, intercambio i\'onico \\
& pH (1~\% p/v) & 4,8 & 5,2 & --- \\
& SO$_2$ & 100 & 250 & ppm \\
& Viscosidad a 60~\textdegree C & --- & 5.000 & cP (m\'aximo especificado) \\
& Color (densidad \'optica) & --- & 1,0 & --- \\
\bottomrule
\end{tabularx}
\end{table}

\subsection{Condiciones Ambientales del Sitio}

La planta de Ingredion~S.A. se ubica en la zona industrial de Cali, Valle del Cauca, Colombia. Las condiciones ambientales de dise\~no se resumen en la Tabla~\ref{tab:environmental_conditions}.

\begin{table}[H]
\centering
\caption{Condiciones ambientales del sitio de instalaci\'on (Cali, Colombia).}
\label{tab:environmental_conditions}
\setcounter{itemcount}{0}
\begin{tabularx}{\textwidth}{@{}NXclc@{}}
\toprule
\multicolumn{1}{c}{\textbf{\'Item}} & \textbf{Par\'ametro Ambiental} & \textbf{Valor} & \textbf{Unidad} & \textbf{Referencia} \\
\midrule
& Temperatura ambiente m\'axima & 38,0 & \textdegree C & IDEAM 2023 \\
& Temperatura ambiente m\'inima & 15,0 & \textdegree C & IDEAM 2023 \\
& Temperatura ambiente promedio & 26,5 & \textdegree C & IDEAM 2023 \\
& Humedad relativa m\'axima & 85,0 & \% & IDEAM 2023 \\
& Humedad relativa m\'inima & 45,0 & \% & IDEAM 2023 \\
& Presi\'on atmosf\'erica & 101,3 & kPa & NOAA 2022 \\
& Altitud sobre el nivel del mar & 960 & msnm & IGAC 2021 \\
& Precipitaci\'on m\'axima anual & 1200,0 & mm & IDEAM 2023 \\
\bottomrule
\end{tabularx}
\end{table}

\subsection{Geometr\'ia del Tanque de Almacenamiento}

El tanque de almacenamiento de glucosa es un recipiente vertical cil\'indrico con fondo toriesf\'erico y tapa c\'onica, fabricado en acero inoxidable SS316L. El fondo toriesf\'erico constituye el elemento principal de este estudio porque aloja la chaqueta de media ca\~na. La Tabla~\ref{tab:tank_geometry} resume las dimensiones principales.

\begin{table}[H]
\centering
\caption{Dimensiones principales del tanque de almacenamiento de glucosa.}
\label{tab:tank_geometry}
\setcounter{itemcount}{0}
\begin{tabularx}{\textwidth}{@{}NXclc@{}}
\toprule
\multicolumn{1}{c}{\textbf{\'Item}} & \textbf{Par\'ametro Geom\'etrico} & \textbf{Valor} & \textbf{Unidad} & \textbf{Componente} \\
\midrule
& Di\'ametro exterior del cuerpo (OD) & 5.276 & mm & Cuerpo cil\'indrico \\
& Di\'ametro interior del cuerpo (ID) & 5.264 & mm & Cuerpo cil\'indrico \\
& Espesor del cuerpo cil\'indrico & 6,0 & mm & Cuerpo (existente) \\
& Altura de la secci\'on cil\'indrica & 9.670 & mm & Cuerpo cil\'indrico \\
& Di\'ametro exterior del fondo (OD) & 5.282 & mm & Fondo toriesf\'erico \\
& Espesor del fondo toriesf\'erico & 9,0 & mm & Fondo (construido) \\
& Altura del fondo toriesf\'erico & 1.266 & mm & Fondo toriesf\'erico \\
& Radio de corona $L$ (ASME F\&D) & 5.264 & mm & Fondo toriesf\'erico \\
& Radio de nudillo $r$ (ASME F\&D) & 315,8 & mm & Fondo toriesf\'erico \\
\midrule
& \textbf{Capacidad nominal (100\%)} & \textbf{222,8} & \textbf{m\textsuperscript{3}} & \textbf{Cilindro + Fondo} \\
& Volumen al 80\% de llenado & 178,2 & m\textsuperscript{3} & Operaci\'on normal \\
\bottomrule
\end{tabularx}
\end{table}

La capacidad nominal se define t\'ecnicamente como el volumen contenido desde la base del fondo toriesf\'erico hasta el nivel superior del cuerpo cil\'indrico. El volumen de la tapa c\'onica superior se reserva como espacio de vapor. La Figura~\ref{fig:fondo_toriesferico} presenta el plano de elevaci\'on del fondo toriesf\'erico con las dimensiones principales, el detalle de la transici\'on cuerpo--fondo y la configuraci\'on hidr\'aulica de la chaqueta. En el plano se identifican dos entradas de agua caliente y dos salidas de agua de retorno, distribuidas de forma sim\'etrica para mejorar la uniformidad t\'ermica del fondo.

\begin{figure}[H]
\centering
\includegraphics[width=0.85\textwidth]{../../Data/planos/fondo.png}
\caption{Plano de elevaci\'on del fondo toriesf\'erico del tanque de glucosa. Se indican las dimensiones principales del radio de corona, radio de nudillo, la transici\'on con el cuerpo cil\'indrico y la configuraci\'on de doble entrada y doble salida de la chaqueta de media ca\~na.}
\label{fig:fondo_toriesferico}
\end{figure}

El sistema de calentamiento consiste en una chaqueta de media ca\~na de perfil rectangular soldada en espiral sobre la superficie exterior del fondo toriesf\'erico. Las caracter\'isticas de la media ca\~na se presentan en la Tabla~\ref{tab:halfpipe_geometry}.

\begin{table}[H]
\centering
\caption{Geometr\'ia de la chaqueta de media ca\~na rectangular en espiral.}
\label{tab:halfpipe_geometry}
\setcounter{itemcount}{0}
\begin{tabularx}{\textwidth}{@{}NXclc@{}}
\toprule
\multicolumn{1}{c}{\textbf{\'Item}} & \textbf{Par\'ametro} & \textbf{Valor} & \textbf{Unidad} & \textbf{Observaci\'on} \\
\midrule
& Ancho interno del perfil rectangular & 141,0 & mm & Lado mayor \\
& Altura interna del perfil rectangular & 45,5 & mm & Lado menor \\
& Espesor de l\'amina del perfil & 4,5 & mm & Acero SS316L \\
& Paso de la espiral & 196 & mm & Centro a centro \\
& \'Area de contacto con el tanque & 14,0 & m\textsuperscript{2} & Superficie efectiva \\
& Entrada central de agua caliente & 45 & \% & Zona central del \'area total \\
& Entrada perif\'erica de agua caliente & 55 & \% & Zona perif\'erica del \'area total \\
& Salida central de agua de retorno & -- & -- & Zona central del \'area total \\
& Salida perif\'erica de agua de retorno & -- & -- & Zona perif\'erica del \'area total \\
& Di\'ametro hidr\'aulico $D_h$ & 68,8 & mm & $4A_c/P_m$ \\
& \'Area de secci\'on transversal $A_c$ & 64,16 & cm\textsuperscript{2} & $141 \times 45,5$~mm \\
& Espesor de aislamiento recomendado & 50,8 & mm & Lana mineral (2''), calculado \\
& Tuber\'ia de alimentaci\'on de agua & 2'' & Sch 40 & Entrada/salida \\
\bottomrule
\end{tabularx}
\end{table}

La Figura~\ref{fig:espiral_mediacana} muestra la vista en planta de la distribuci\'on en espiral de la media ca\~na sobre el fondo toriesf\'erico, incluyendo las dos entradas y dos salidas de agua caliente y el detalle del paso entre espiras.

\begin{figure}[H]
\centering
\includegraphics[width=0.85\textwidth]{../../Data/planos/espiral.png}
\caption{Vista en planta de la chaqueta de media ca\~na en espiral sobre el fondo toriesf\'erico. Se indica la configuraci\'on de doble entrada de agua caliente (45~\% en zona central y 55~\% en zona perif\'erica) y doble salida de agua de retorno, junto con el paso de la espiral de 196~mm. La distribuci\'on sim\'etrica de las conexiones reduce los gradientes t\'ermicos entre espiras interiores y exteriores.}
\label{fig:espiral_mediacana}
\end{figure}

\subsection{Propiedades Fisicoqu\'imicas de la Glucosa Globe 1130}

El fluido almacenado es glucosa Globe 1130, un jarabe de glucosa de alta concentraci\'on. Seg\'un la ficha t\'ecnica oficial del fabricante (Globe\textsuperscript{\textregistered}~1130 Corn Syrup/Glucose~011420, Ingredion~2017), se trata de un jarabe de ma\'iz de conversi\'on regular, intercambio i\'onico, con equivalente de dextrosa (DE) de 40,5--46,5 y concentraci\'on de s\'olidos seca de 79,7--81,5~\%. La Tabla~\ref{tab:ingredion_spec} resume las especificaciones oficiales del producto.

\begin{table}[H]
\centering
\caption{Especificaciones oficiales del jarabe Globe\textsuperscript{\textregistered}~1130 (ficha t\'ecnica Ingredion~011420, 2017).}
\label{tab:ingredion_spec}
\setcounter{itemcount}{0}
\begin{tabularx}{\textwidth}{@{}NlccX@{}}
\toprule
\multicolumn{1}{c}{\textbf{\'Item}} & \textbf{Propiedad} & \textbf{M\'in} & \textbf{M\'ax} & \textbf{Observaci\'on} \\
\midrule
& Sustancia seca [\%] & 79,7 & 81,5 & \textdegree Brix equivalente \\
& Equivalente de dextrosa (DE) & 40,5 & 46,5 & Conversi\'on regular \\
& Color, CP & -- & 1,2 & -- \\
& SO\textsubscript{2} [ppm] & -- & $<$ 10 & -- \\
& Cenizas [\%] & -- & 0,05 & -- \\
& Almid\'on aparente & \multicolumn{2}{c}{Negativo} & -- \\
\midrule
\multicolumn{5}{@{}l}{\textit{Perfil de carbohidratos (\% base seca, t\'ipico):}} \\
& Dextrosa & \multicolumn{2}{c}{21\%} & -- \\
& Maltosa & \multicolumn{2}{c}{13\%} & -- \\
& Maltotriosa & \multicolumn{2}{c}{11\%} & -- \\
& Sac\'aridos superiores (DP4+) & \multicolumn{2}{c}{55\%} & -- \\
\midrule
\multicolumn{5}{@{}l}{\textit{Densidad y viscosidad vs.\ temperatura (ficha t\'ecnica):}} \\
& 26,7~\textdegree C (80~\textdegree F) & \multicolumn{2}{c}{$\rho = 1,421$~kg/L, $\mu = 74{.}000$~cP} & -- \\
& 37,8~\textdegree C (100~\textdegree F) & \multicolumn{2}{c}{$\rho = 1,415$~kg/L, $\mu = 17{.}400$~cP} & -- \\
& 48,9~\textdegree C (120~\textdegree F) & \multicolumn{2}{c}{$\rho = 1,409$~kg/L, $\mu = 5{.}400$~cP} & -- \\
\midrule
& Temperatura de almacenamiento recomendada & 49 & 54 & \textdegree C (120--130~\textdegree F) \\
& Certificaciones & \multicolumn{2}{c}{Kosher pareve, Halal} & FCC, 21~CFR~168.120 \\
\bottomrule
\end{tabularx}
\end{table}

La concentraci\'on de dise\~no se establece en 80,6~\textdegree Brix, valor nominal medio de la ficha t\'ecnica. A temperaturas inferiores a 40~\textdegree C la glucosa Globe 1130 es extremadamente viscosa y se comporta pr\'acticamente como un fluido semis\'olido; a temperaturas superiores a 50~\textdegree C su viscosidad disminuye significativamente y el fluido puede ser bombeado y descargado con facilidad. Esta dependencia de la viscosidad con la temperatura es el factor determinante en el dise\~no del sistema de calentamiento.

La Tabla~\ref{tab:glucose_properties} presenta las propiedades termof\'isicas de la glucosa Globe 1130 a 80,6~\textdegree Brix en funci\'on de la temperatura. La viscosidad se obtuvo mediante un modelo de Vogel-Fulcher-Tammann (VFT) calibrado a los tres puntos de la ficha t\'ecnica; la densidad proviene del ajuste lineal a los mismos tres puntos. El calor espec\'ifico y la conductividad t\'ermica se calcularon con las correlaciones de Choi y Okos para soluciones de carbohidratos.

\begin{table}[H]
\centering
\caption{Propiedades termof\'isicas de la glucosa Globe 1130 ($\sim$80,6~\textdegree Brix) en funci\'on de la temperatura.}
\label{tab:glucose_properties}
\setcounter{itemcount}{0}
\begin{tabularx}{\textwidth}{@{}NcccccX@{}}
\toprule
\multicolumn{1}{c}{\textbf{\'Item}} & \textbf{$T$ [\textdegree C]} & \textbf{$\rho$ [kg/m\textsuperscript{3}]} & \textbf{$\mu$ [cP]} & \textbf{$C_p$ [J/(kg$\cdot$\textdegree C)]} & \textbf{$k$ [W/(m$\cdot$\textdegree C)]} & \textbf{Fuente} \\
\midrule
& 25 & 1.421,6 & 212{.}000 & 2.092 & 0,304 & Modelo VFT \\
& 26,7 & 1.421,0 & 74{.}000 & 2.095 & 0,306 & Ficha t\'ecnica \\
& 37,8 & 1.415,0 & 17{.}400 & 2.106 & 0,316 & Ficha t\'ecnica \\
& 48,9 & 1.409,0 & 5{.}400 & 2.117 & 0,327 & Ficha t\'ecnica \\
& 57 & 1.404,6 & 2{.}600 & 2.131 & 0,340 & Modelo VFT \\
& 60 & 1.403,0 & 2{.}060 & 2.134 & 0,344 & Modelo VFT \\
& 70 & 1.397,6 & 986 & 2.143 & 0,352 & Modelo VFT \\
& 80 & 1.392,2 & 522 & 2.151 & 0,359 & Modelo VFT \\
\bottomrule
\end{tabularx}
\end{table}

La caracter\'istica m\'as relevante de la glucosa Globe 1130 para este estudio es la ca\'ida dram\'atica de la viscosidad con la temperatura: de 212{.}000~cP a 25~\textdegree C a 2{.}060~cP a 60~\textdegree C, lo que representa una reducci\'on de dos \'ordenes de magnitud. Este comportamiento se modela mediante una ecuaci\'on de Vogel-Fulcher-Tammann (VFT):
\begin{equation}
\ln(\mu) = A + \frac{B}{T + C}
\label{eq:vft}
\end{equation}
donde $A = -9,350$, $B = 1{.}278,0$~\textdegree C y $C = 66,90$~\textdegree C, con $\mu$ en Pa$\cdot$s y $T$ en \textdegree C. Los tres par\'ametros se determinaron mediante ajuste exacto a los tres puntos de densidad y viscosidad de la ficha t\'ecnica Ingredion~011420.

\subsection{Propiedades del Agua de Calentamiento}

El fluido de servicio en la chaqueta de media ca\~na es agua caliente proveniente del sistema de servicios industriales de la planta. Las propiedades del agua se eval\'uan a la temperatura media de operaci\'on en cada escenario mediante correlaciones simplificadas validadas en el rango de 40 a 80~\textdegree C. La Tabla~\ref{tab:water_properties} presenta valores representativos.

\begin{table}[H]
\centering
\caption{Propiedades termof\'isicas del agua de calentamiento en el rango de operaci\'on.}
\label{tab:water_properties}
\setcounter{itemcount}{0}
\begin{tabularx}{\textwidth}{@{}NcccccX@{}}
\toprule
\multicolumn{1}{c}{\textbf{\'Item}} & \textbf{$T$ [\textdegree C]} & \textbf{$\rho$ [kg/m\textsuperscript{3}]} & \textbf{$\mu$ [cP]} & \textbf{$C_p$ [J/(kg$\cdot$\textdegree C)]} & \textbf{$k$ [W/(m$\cdot$\textdegree C)]} & \textbf{Pr} \\
\midrule
& 50 & 989,6 & 0,547 & 4.167 & 0,641 & 3,56 \\
& 60 & 985,3 & 0,467 & 4.164 & 0,650 & 2,99 \\
& 65 & 982,9 & 0,433 & 4.163 & 0,654 & 2,75 \\
& 70 & 980,3 & 0,404 & 4.161 & 0,657 & 2,56 \\
& 75 & 977,5 & 0,378 & 4.160 & 0,660 & 2,38 \\
& 80 & 974,4 & 0,355 & 4.158 & 0,662 & 2,23 \\
\bottomrule
\end{tabularx}
\end{table}

\subsection{Condiciones de Operaci\'on y Escenarios de An\'alisis}

El tanque de glucosa opera a presi\'on atmosf\'erica sin agitaci\'on mec\'anica. El calentamiento se realiza exclusivamente por medio de agua caliente que circula por la chaqueta de media ca\~na en espiral sobre el fondo toriesf\'erico. La glucosa se recibe a temperaturas entre 55 y 60~\textdegree C y se despacha a carrotanque en el rango de 57 a 60~\textdegree C, necesario para reducir la viscosidad a niveles que permitan el bombeo.

Se definieron cinco escenarios de an\'alisis para evaluar el desempe\~no t\'ermico del sistema, los cuales se resumen en la Tabla~\ref{tab:escenarios}.

\begin{table}[H]
\centering
\caption{Escenarios de an\'alisis t\'ermico definidos para el estudio.}
\label{tab:escenarios}
\setcounter{itemcount}{0}
\begin{tabularx}{\textwidth}{@{}NlccccX@{}}
\toprule
\multicolumn{1}{c}{\textbf{\'Item}} & \textbf{Escenario} & \textbf{$V_g$ [m\textsuperscript{3}]} & \textbf{$T_{w,in}$ [\textdegree C]} & \textbf{$v_w$ [m/s]} & \textbf{$Q_w$ [m\textsuperscript{3}/h]} & \textbf{Observaci\'on} \\
\midrule
& Escenario 1 & 24,0 & 65 & 1,34 & 30,9 & Descarga continua con alimentaci\'on simult\'anea \\
& Escenario 2 & 178,2 & 65 & 2,50 & 57,7 & Calentamiento desde 25~\textdegree C, tanque al 80\% \\
& Escenario 3 & 178,2 & 75 & 2,50 & 57,7 & Calentamiento desde 25~\textdegree C, tanque al 80\% \\
& Escenario 4 & 136,9 & 75 & 2,50 & 57,7 & Capacidad operativa diaria: 5 descargas de 24~ton \\
& Escenario 5 & 178,2 & 75 & 2,50 & 57,7 & Ciclo de 5 descargas con recalentamiento entre descargas \\
\bottomrule
\end{tabularx}
\end{table}

El Escenario~1 eval\'ua el calentamiento de un volumen parcial de 24~m\textsuperscript{3} de glucosa con un caudal de agua de 30,9~m\textsuperscript{3}/h, incluyendo la descarga de 24~toneladas a carrotanque al alcanzar la temperatura objetivo. Los Escenarios~2 y 3 representan la condici\'on de arranque desde fr\'io: el tanque al 80~\% de su capacidad (178,2~m\textsuperscript{3}, equivalentes a 253~toneladas de glucosa), con velocidad del agua en la media ca\~na de 2,5~m/s. La diferencia entre ambos escenarios es la temperatura del agua de entrada (65~\textdegree C frente a 75~\textdegree C). El Escenario~4 modela la capacidad operativa diaria del sistema: cinco descargas consecutivas de 24~toneladas dentro de un ciclo de 24~horas, con flujo medio de 5.000~kg/h. El Escenario~5 simula el ciclo din\'amico completo de cinco descargas con recalentamiento entre descargas, partiendo de un inventario caliente.
